\begin{abstract}
A run-to-run standard error of a benchmark average that sums per-benchmark variances ignores covariance between benchmarks, and in simulation a margin comparison using it declares a difference in 0.113 of null replicates, against a nominal 0.05. SNAP (Seed Noise Across Phenotypes) estimates that covariance from cross-half products of run-centred item-half scores on a fixed battery of likelihood-scored multiple-choice benchmarks. On 375 DataDecide runs at each configuration's last shared checkpoint, the ten-benchmark per-byte margin average has inflation, the ratio of its standard deviation to that under independence, of 1.244 \ci{1.143}{1.338}, or 1.237 \ci{1.137}{1.331} from replicate standard deviations alone, which need no item split. BoolQ carries 68\% of the covariance trace, and single-benchmark removals give 1.096 to 1.786. Acting alone, a uniform item component shared by all benchmarks could carry at most about a third of the margin excess before inflation across scoring formats leaves \ci{0.801}{1.035}. The accuracy interval \ci{0.993}{1.157} includes one, at 0.135 power against a true 1.10, and on DataDecide's recipe comparisons the correction removes at most 12 of the 173 to 223 calls on pairs with a resolved 1B ordering. A pre-specified test on four held-out tasks fails its lower-limit rule at 1.217 \ci{0.872}{1.498}, against 1.240 \ci{1.054}{1.397} for the original ten. \outcome{\Ronecase}{1}{R1 pass}{\outcome{\Rtwocase}{1}{R2 not supported}{Under rules fixed before scoring, nine seeds at five PolyPythias sizes give margin inflation of \pending{R1 lambda} with interval \ci{\pending{lo}}{\pending{hi}}, and pairing that battery with a disjoint item bank leaves \pending{cross lambda}, which gives no support to the shared-item explanation.}\outcome{\Rtwocase}{2}{R2 supported}{Under rules fixed before scoring, nine seeds at five PolyPythias sizes give margin inflation of \pending{R1 lambda} with interval \ci{\pending{lo}}{\pending{hi}}, although a disjoint item bank supports a shared item component.}\outcome{\Rtwocase}{3}{R2 undecided}{Under rules fixed before scoring, nine seeds at five PolyPythias sizes give margin inflation of \pending{R1 lambda} with interval \ci{\pending{lo}}{\pending{hi}}, although a disjoint item bank leaves the source of that inflation undecided.}}\outcome{\Ronecase}{2}{R1 fail}{On a second family, nine seeds at five PolyPythias sizes give margin inflation of \pending{R1 lambda} with interval \ci{\pending{lo}}{\pending{hi}}, which fails a rule fixed before scoring at a design with 40 within-configuration contrasts, and a disjoint item bank \outcome{\Rtwocase}{1}{R2 not supported}{leaves \pending{cross lambda}, which gives no support to the shared-item explanation}\outcome{\Rtwocase}{2}{R2 supported}{supports a shared item component}\outcome{\Rtwocase}{3}{R2 undecided}{leaves the source of that inflation undecided}.} Beyond the pre-specified test and two PolyPythias rules, analyses are exploratory. We recommend reporting the replicate standard deviation of the battery average with a benchmark-removal check.
\end{abstract}

\section{Introduction}
A benchmark average can change across training runs even when the evaluation items stay fixed. For an equally weighted average of $K$ benchmarks, independent run deviations give variance $K^{-2}\operatorname{tr}\Sig$, where $\Sig$ is their cross-benchmark covariance matrix. Assuming independence drops the off-diagonal terms and weakens a comparison's error control. In simulation with margin-like covariance, such a comparison of two independently trained configurations declares a difference in 0.113 of replicates with a true gap of zero, against a nominal 0.05 (Appendix~\ref{app:supplementary}). On DataDecide's recipe comparisons the correction changes few calls (Section~\ref{sec:prediction}), so the understatement shows mainly in the error bar of one model's battery average and in planning replicate runs.

We estimate item-general covariance from two disjoint item halves per benchmark, centring both half scores across each configuration's replicate runs and taking their cross-half products. Under zero-mean item noise and conditional independence across halves, they estimate run covariance, including its diagonal. We adapt the repeated-measurement construction of quantitative genetics and call it SNAP, Seed Noise Across Phenotypes (Appendix~\ref{app:derivations}).

Our empirical study uses DataDecide \citep{magnusson2025}, with 25 data recipes and five of its model sizes, each with three released replicates, and those 125 configurations supply 375 runs across ten benchmarks, scored at the largest checkpoint step each configuration's replicates share. Our primary analysis reads released item scores without model inference and compares per-byte log-likelihood margins with accuracy. Margins carry more seed signal, with per-benchmark reliability averaging 0.652 against 0.324 for accuracy (run variance's share of a benchmark's replicate variance; Section~\ref{sec:composition}), so we treat margins as the primary scale and accuracy as a bound.

Margin inflation $\Lam$, the ratio of the average's standard deviation to its independent value, is 1.244 \ci{1.143}{1.338}, while accuracy gives 1.078 \ci{0.993}{1.157} at the scored checkpoint (Table~\ref{tab:primary}).

Because we built our estimator on these ten benchmarks, we pre-specified a test on SciQ \citep{welbl2017sciq}, MedMCQA \citep{pal2022medmcqa}, and DROP \citep{dua2019drop} and CoQA \citep{reddy2019coqa} in multiple-choice form, on the same checkpoints at 530M, 750M, and 1B. Its rule requires the lower interval limit of held-out margin inflation to exceed one. The test fails at its full scope of 225 runs, where held-out inflation is 1.217 \ci{0.872}{1.498} against 1.240 \ci{1.054}{1.397} for the original ten on the same runs (Section~\ref{sec:heldout}).

A single model family whose replicates mix default and auxiliary runs on different schedules leaves open whether the inflation belongs to DataDecide's design. We therefore score all 45 PolyPythias runs at five sizes, where each of nine seeds reruns one recipe, and pair that battery with a disjoint bank of 6,808 items built with the OLMES harness \citep{gu2025olmes} at a pinned commit. Two rules fixed before scoring read the results, and in simulation the first passes in 0.259 of replicates at a true 1.244 and the second returns undecided in 0.761, so a first-rule pass confirms the effect, a failure bounds it or counts against transfer, and an undecided second rule leaves the shared-item question open. Section~\ref{sec:family} reports that margin inflation there is \pending{R1 lambda} with interval \ci{\pending{lo}}{\pending{hi}} and that the shared-item explanation is \outcome{\Rtwocase}{1}{R2 not supported}{not supported}\outcome{\Rtwocase}{2}{R2 supported}{supported}\outcome{\Rtwocase}{3}{R2 undecided}{undecided}.

This paper makes four contributions to measuring run covariance. (i) We define the item-general run covariance that cross-half products estimate, identified up to an item component shared across benchmarks and halves. (ii) The replicate standard deviation of the battery average gives the uncertainty of rerunning the released items and is what we recommend reporting, while the split adds per-benchmark run variances net of item noise and a redrawn-bank target, untested on DataDecide, and the two part where item noise is large, as for accuracy without BoolQ at 1.285 against 1.558 (Section~\ref{sec:estimator}). (iii) Algorithm~\ref{alg:snap} gives the measurement recipe, with near-nominal intervals unless a run effect follows the size band. (iv) We measure it across the 375 released runs, where every recipe, recipe-pair and benchmark deletion, along with four of the five size-band deletions, leaves a margin interval that excludes one.
